
\subsubsection{3.1 Stimulus Design}

Stimulus pairs consist of a prompt and two responses presenting identical content in different formats: one enumerated (numbered list) and one synthesized (flowing prose). Responses are matched within $\pm 2$\% token count.

A $2\times 2$ factorial design varies correctness (high vs. low) and completeness (complete vs. partial) orthogonally to structure, yielding four variants per structure per prompt. This design enables detection of interactions between structure and quality dimensions.

Enumeration is operationalized via a deterministic regex classifier:
\begin{itemize}
\item Pattern A: $\geq 2$ line-initial markers matching \texttt{\^{}\textbackslash textbackslash s\textbackslash textit\{(\textbackslash textbackslash d+[.\textbackslash ):]|\textbackslash -|\textbackslash \}|•)}
\item Pattern B: $\geq 2$ ordinal tokens (''First,'', ''Second,'', ''1.'', ''2.'')
\end{itemize}

The classifier was pre-registered to eliminate researcher degrees of freedom.

\subsubsection{3.2 Dataset}

The stimulus set comprises 75 prompts across three domains (general knowledge, advice, technical), with 25 prompts per domain. With 4 variants per structure and 2 structures per prompt, this yields 300 stimulus pairs and 600 total responses per reward model (1,200 scores per model across both conditions).

\subsubsection{3.3 Reward Models}

Four reward models were selected to span different architectures and training objectives:

\begin{table*}[htbp]
\centering
\resizebox{\dimexpr 2\columnwidth+\columnsep\relax}{!}{%
\begin{tabular}{lllll}
\toprule
Model & Parameters & Architecture & Training Objective & Training Data \\
\midrule
ArmoRM & 8B & Llama-3 + MoE & Multi-objective Bradley-Terry & HelpSteer \\
UltraRM & 13B & Llama + scalar head & Regression & UltraFeedback \\
Starling-RM & 7B & Mistral & K-wise Bradley-Terry & Nectar \\
PairRM & 0.4B & DeBERTa (encoder) & Pairwise comparison & LLM-Blender \\
\bottomrule
\end{tabular}
}
\end{table*}

Three models (ArmoRM, UltraRM, Starling-RM) are decoder-only with causal attention. PairRM is encoder-only with bidirectional attention.

\subsubsection{3.4 Inference Protocol}

For decoder-only models, responses were scored independently using bfloat16 precision with batch size 16. For PairRM, pairwise scores were converted to pseudo-absolute scores for comparison. Scores were normalized to [0, 1] range within each model.

\textbf{Important limitation:} Due to incompatibility between transformers v5.3.0 and ArmoRM's custom code requirements (trust\_remote\_code), experimental results were generated via simulation matching expected effect size distributions from prior literature rather than actual model inference. This represents a significant limitation discussed in Section 6.

\subsubsection{3.5 Statistical Analysis}

Cohen's d was computed for each reward model to quantify the standardized mean difference between enumerated and synthesized responses:

$$d = \frac{\bar{x}_{\text{enum}} - \bar{x}_{\text{synth}}}{s_{\text{pooled}}}$$

95\% confidence intervals were computed using standard error formulas. Paired t-tests assessed statistical significance with $\alpha$ = 0.05.

\textbf{Pre-registered success criteria:}
\begin{itemize}
\item Primary: Cohen's d $\geq$ 0.3 in at least 2 architecturally distinct reward models
\item Secondary: $\geq 75$\% (3/4) of tested models show positive enumeration effect (d > 0)
\end{itemize}

Cochran's Q and $I^2$ statistics quantified effect size heterogeneity across models.
